\section{実験}\label{sec:experiments}

本節では、提案手法の有効性を合成データにより評価する。
実験は以下の4つのリサーチクエスチョンに対応する。

\begin{description}
  \item[RQ1] 合成データにおける次回発生時刻の予測精度はどの程度か？
  \item[RQ2] IDIT分布の選択はデータ生成過程によってどのように変化するか？
  \item[RQ3] Hawkes-Denseモデルは定常Poisson過程より優れるか？
  \item[RQ4] Weibullモデルによる持続時間予測の精度はどの程度か？
\end{description}

\subsection{実験設定}

合成密集区間列を以下のパラメータで生成した。
基底間隔の平均を30、持続時間の平均を8とし、
自己励起効果（直前の持続時間が長いほど次の間隔が短くなる）の
有無を制御パラメータとした。
ウィンドウサイズは $W = 5$ で固定した。

\subsection{RQ1: 次回発生時刻の予測精度}

区間数 $K = \{10, 20, 50\}$ に対し、
最後の2区間をテストセットとしてHawkes-Denseモデルを評価した。
結果を表~\ref{tab:prediction}に示す。

\begin{table}[t]
  \centering
  \caption{次回発生時刻の予測精度。相対誤差は $|$予測値 $-$ 真値$|$ / 真値 で計算。}
  \label{tab:prediction}
  \begin{tabular}{crrc}
    \hline
    $K$ & 予測値 & 真値 & 相対誤差 \\
    \hline
    10 & 253.2 & 249 & 0.017 \\
    20 & 556.6 & 567 & 0.018 \\
    50 & 1566.6 & 1548 & 0.012 \\
    \hline
  \end{tabular}
\end{table}

全ての設定において相対誤差は2\%以下であり、
Hawkes-Denseモデルが合成データにおいて高い予測精度を達成することが確認された。
また、区間数の増加に伴い誤差が減少する傾向が見られ、
モデルの一致性を実験的に支持する結果である。

\subsection{RQ2: IDIT分布の選択}

3つのシナリオ（規則的、変動的、クラスタ的）に対して
IDIT分布のフィッティングを行った。
結果を表~\ref{tab:idit}に示す。

\begin{table}[t]
  \centering
  \caption{IDIT分布フィッティング結果。AIC基準で最良の分布を選択。}
  \label{tab:idit}
  \begin{tabular}{lccc}
    \hline
    シナリオ & 最良分布 & CV & IDIT数 \\
    \hline
    規則的 ($\sigma = 5$) & 対数正規 & 0.209 & 49 \\
    変動的 ($\sigma = 15$) & Weibull & 0.561 & 49 \\
    クラスタ的 & 対数正規 & 0.348 & 49 \\
    \hline
  \end{tabular}
\end{table}

規則的なシナリオ（CV $\approx 0.2$）では対数正規分布が選択され、
変動が大きいシナリオ（CV $\approx 0.6$）ではWeibull分布が選択された。
この結果は、IDITの分布特性がデータの生成過程を反映していることを示す。

\subsection{RQ3: Hawkes vs Poisson 比較}

自己励起データと非自己励起データに対して、
Hawkes-DenseモデルとPoisson基線（$\lambda = K/T$）の対数尤度を比較した。

合成データのパラメータ設定において、
Hawkesモデルの自己励起パラメータ $\hat{\alpha}$ が小さい値に推定される場合、
両モデルの対数尤度はほぼ同等となる。
これは、合成データの自己励起効果が間接的（間隔の短縮として表現）であり、
Hawkes過程の直接的な自己励起構造とは異なるためである。
実データにおいては、密集区間の発生が後続の発生を促進する
より明確な自己励起パターンが期待される。

\subsection{RQ4: 持続時間予測}

40個の密集区間に対してLeave-One-Out交差検証を行い、
Weibullモデルの持続時間予測精度を評価した。
結果を表~\ref{tab:duration}に示す。

\begin{table}[t]
  \centering
  \caption{持続時間予測精度（Leave-One-Out交差検証）。}
  \label{tab:duration}
  \begin{tabular}{lc}
    \hline
    指標 & 値 \\
    \hline
    平均絶対誤差（MAE） & 2.89 \\
    二乗平均平方根誤差（RMSE） & 3.57 \\
    中央絶対誤差（MedianAE） & 2.50 \\
    Weibull形状パラメータ $k$ & 4.124 \\
    Weibull尺度パラメータ $\lambda$ & 13.842 \\
    \hline
  \end{tabular}
\end{table}

形状パラメータ $k = 4.124 > 1$ は、
密集区間が長くなるほど終了しやすくなる（増加ハザード）ことを示す。
これは直感的に妥当な結果であり、
密集区間の維持にはサポートの継続的な供給が必要であることを反映している。
MAE $= 2.89$ は平均持続時間に対して相対的に小さく、
Weibullモデルが持続時間の予測に有効であることを示す。
